\scanpage{63}
Il ne semble pas que $5^\circ$ et $6^\circ$ s'expriment de façon simple en termes des $M_n$. Mais on en profite pour reformuler les données ainsi. Pour tt $n\geq0$, soit
\[
\mathcal F_n\subset\operatorname{Hom}_{\mathcal K}(I^n,I)\qquad\deleted{\unsure{\alpha\in\mathcal F}}
\]
l'ens $\operatorname{Hom}_{\mathcal M}(I_0^n,I_0)$. Grâce à $6^\circ$, la connaissance des $\mathcal F_n$ implique celle des $M_n$. Quelles conditions à mettre ?

\deleted{(F1)} Posons $\mathcal F_{\unsure{m,n|n,m}}\subset\operatorname{Hom}_{\mathcal K}(I^n,I^m)$ \deleted{défini} (morphismes $\mathcal F$-modérés) par $\varphi=(\varphi_1-\varphi_m)$ $\mathcal F$-modéré ssi les $\varphi_i\in\mathcal F_n$, on \deleted{\unsure{aura}} aura

\inserted{$(F0)$ $\mathcal F_n$ contient les $\operatorname{pr}_i:I^n-I$.}

\deleted{\unsure{F1}} $(F2)$ Stabilité par compositions des $\mathcal F_{n,m}$.

On pose $M_n=$ ens des sous-objets de $I^n$ \inserted{intersections finies de ss-objets} de la forme $\varphi^{-1}(\alpha)$, où $\varphi\in\mathcal F_n$, $\alpha\in\mathcal F_0:e\xrightarrow{\alpha}I$.
Alors $M_n$ satisfait $M1,M2$, pour avoir $M3$ bis (forme forte) on doit admettre

$(F3)$ Si $\varphi_1-\varphi_m\in\mathcal F_n$ d'où $I^n\longrightarrow I^m$

\inserted{$\Updownarrow$ $(M3\text{bis})$}
et si \unsure{---}

$(F4)$ a) $\exists\,0:e\longrightarrow I$ \inserted{$\in\mathcal F_0$} tel que $\forall\alpha\in\mathcal F_1$, $\alpha:e-I$, on ait
\[
\operatorname{Im}\alpha=\bigcap_{1\leq i\leq m}\varphi_i^{-1}(0),\quad
\text{pour }\varphi_1-\varphi_m\in\mathcal F_1
\]
\inserted{$\in M_1$ i.e.}

b) $\operatorname{Sup}(e\times I;I\times e)\subset I^2$ est $\in M_2$, i.e. est de la forme $\bigcap_i\varphi_i^{-1}(0)$, où $\varphi_i\in\mathcal F_2$.

c) $\exists\,\varepsilon_1,\varepsilon_2\in\mathcal F_0$, $\varepsilon_1,\varepsilon_2:e-I$, tel que $e\amalg e-I$ soit mono.

Il \deleted{\unsure{faut}} exprimer aussi que pour \deleted{(5) Soit} $f:\deleted{X}\inserted{I^n}\to I$ \deleted{$\unsure{\mathcal F_n}$} \inserted{$\unsure{X}\subset I^n$, se prolonge en $\unsure{\bar f}$} \deleted{d'où} $\Gamma_f\subset I^{n+1}$,

Pour que $\deleted{\unsure{\bar f}}\ f\in\mathcal F_n$, il f. et s. que $\Gamma_f\in M_{n+1}$;
\deleted{i.e. a) Si $f\in\mathcal F_n$, $\exists$}
le ``il faut'' équivaut au fait que $\Delta_I\subset I\times I\longrightarrow{}\in M_2$; le ``il suffit'' doit être postulé indépendamment

\margin{64}
